\chapter{Backtests vectorisés}\label{ch:rs:vectorised-backtests}

Un \omterm{def:rs:vectorised-backtests:backtest}{backtest} en une seule ligne de l'inversion à un jour sur le marché synthétique (acheter les perdants de la veille, vendre ses gagnants, chaque jour) affiche un ratio de Sharpe
de 7.2 sur dix ans. Écrite autrement, en multipliant le portefeuille construit à partir du cours de clôture du jour par le rendement du jour, elle affiche
$-65$. Aucun des deux n'est un résultat. La première opère à la clôture qui a servi à calculer le signal, ignore les coûts de transaction sur un portefeuille qui
se renouvelle des trois quarts de lui-même chaque jour, et perdrait de l'argent de l'une quelconque des façons que ce chapitre énumère~; la seconde est un \omterm{def:rs:point-in-time-data-and-the-biases:lookahead}{biais d'anticipation} si
grossier qu'il inverse le signal. Le \omterm{def:rs:vectorised-backtests:backtest}{backtest vectorisé}, quelques opérations sur des tableaux portant sur un panel de poids et de rendements, est l'outil de recherche le plus rapide de la
société, et la plupart du temps le bon~; ce chapitre le construit (\texttt{firm.vecbt}), indique quand il est juste, et
démonte ses mensonges un par un.

\section{Le backtest vectorisé}

\begin{definition}[Backtest, backtest vectorisé]\label{def:rs:vectorised-backtests:backtest}
Un \emph{backtest}\index{backtest} est la simulation d'une règle de trading sur des données historiques, produisant les positions, les transactions, les coûts et
les rendements que la règle aurait eus. Un \emph{backtest vectorisé}\index{backtest vectorisé} les calcule sous forme d'opérations sur des tableaux entiers, sur des panels
(périodes par titres)~: poids cibles, rendements, et coûts en fonction du poids échangé, sans simuler les ordres.
\end{definition}

\begin{definition}[Instant de décision, délai d'exécution, fréquence de rebalancement]\label{def:rs:vectorised-backtests:timing}
L'\emph{instant de décision}\index{instant de décision} d'un portefeuille est le moment où l'information utilisée pour le construire était complète. Le
\emph{délai d'exécution}\index{délai d'exécution} est le nombre de périodes séparant l'instant de décision du premier rendement que le portefeuille rapporte. La
\emph{fréquence de rebalancement}\index{fréquence de rebalancement} est la périodicité à laquelle le portefeuille cible est recalculé et échangé.
\end{definition}

Dans \texttt{firm.vecbt}, un portefeuille décidé à la clôture du jour $t$ avec un délai $L$ est détenu le jour $t + L$~: le délai 0 le multiplie par le rendement qui
l'a produit, le délai 1 suppose qu'il a été échangé à la clôture même d'où il a été calculé, le délai 2 l'échange à la clôture suivante. Le portefeuille dérive avec
les rendements entre deux rebalancements, si bien que le poids échangé lors d'un rebalancement est la cible moins le portefeuille dérivé, et non la cible moins la cible précédente~;
les titres situés hors de l'\omterm{def:rs:universe-symbology-and-corporate-actions:universe}{univers investissable} à l'\omterm{def:rs:vectorised-backtests:timing}{instant de décision} ne sont pas détenus~; un titre qui cesse d'être coté perçoit son \omterm{def:rs:universe-symbology-and-corporate-actions:delisting}{rendement de radiation} et rien
après. Le résultat est un \texttt{BacktestResult} (poids, transactions, rendements bruts, coûts par type, rendements nets, capital), le type que renvoient tous les
backtesteurs de la société et que lit le module de performance du chapitre~22.

\begin{definition}[Rotation du portefeuille, modèle de coûts linéaire]\label{def:rs:vectorised-backtests:turnover}
La \emph{rotation du portefeuille}\index{rotation du portefeuille} (portfolio turnover) d'une période est la moitié de la somme des poids absolus échangés, la fraction du portefeuille
remplacée. Un \emph{modèle de coûts linéaire}\index{modèle de coûts linéaire} (linear cost model) facture à chaque transaction un coût fixe par unité de poids échangé (une demi-fourchette plus
les commissions), indépendant de sa taille.
\end{definition}

\section{Quatre niveaux de fidélité}

\begin{definition}[Niveau de fidélité]\label{def:rs:vectorised-backtests:fidelity}
Le \emph{niveau de fidélité}\index{niveau de fidélité} d'un \omterm{def:rs:vectorised-backtests:backtest}{backtest} est la part du processus de trading qu'il simule. Ce livre en utilise quatre~: le niveau~1,
le \omterm{def:rs:vectorised-backtests:backtest}{backtest vectorisé} sur les rendements de barres, avec des coûts fonctions de la rotation~; le niveau~2, le \omterm{def:rs:vectorised-backtests:backtest}{backtest} événementiel sur barres, avec ordres, exécutions et
une horloge simulée (chapitre~17)~; le niveau~3, le rejeu des messages du carnet d'ordres avec positions dans la file et latence (chapitre~18)~; le niveau~4, le trading en direct
à petite taille ou trading simulé (paper trading) sur données en temps réel (chapitres 19 et 21).
\end{definition}

\begin{omfigure}
\begin{tikzpicture}[font=\small,node distance=3mm,
  lvl/.style={draw=omDef,fill=omDef!8,rounded corners=2pt,align=left,text width=3.2cm,inner sep=4pt,minimum height=3.3cm,anchor=north}]
\node[lvl] (a) at (0,0) {\textbf{Niveau 1}\\ vectorisé\\[2pt] poids $\times$ rendements~;\\ coûts à partir de la rotation~;\\ secondes par passage};
\node[lvl] (b) at (3.7,0) {\textbf{Niveau 2}\\ événementiel, barres\\[2pt] ordres, modèles d'exécution,\\ horloge, calendriers~;\\ minutes par passage};
\node[lvl] (c) at (7.4,0) {\textbf{Niveau 3}\\ rejeu du carnet d'ordres\\[2pt] files, latence,\\ impact propre~;\\ heures par passage};
\node[lvl] (d) at (11.1,0) {\textbf{Niveau 4}\\ en direct ou trading simulé\\[2pt] exécutions réelles,\\ latence réelle~;\\ semaines par réponse};
\draw[-Stealth,thick,omThm] (-1.5,-3.65) -- node[below,font=\footnotesize]{la fidélité et le coût de chaque passage augmentent~; le nombre d'idées testables diminue} (12.6,-3.65);
\end{tikzpicture}

\omcaption{Les quatre niveaux de fidélité des \omterm{def:rs:vectorised-backtests:backtest}{backtests} de la société. Une idée de recherche ne progresse vers la droite que si elle survit à gauche.}\label{fig:rs:vectorised-backtests:levels}
\end{omfigure}

Les niveaux ne sont pas rivaux~: mille idées sont criblées au niveau~1 dans le temps qu'il faut pour en exécuter une au niveau~3, et un résultat de niveau~1 négatif
après coûts n'a pas besoin du niveau~3 pour être rejeté. La question, pour chaque stratégie, est le niveau le plus bas
dont la réponse ne serait pas modifiée par les niveaux supérieurs.

\section{Où le backtest vectorisé est juste}

Il est juste lorsque le trading ne modifie pas la réponse~: signaux lents (périodes de détention de quelques jours à quelques mois), titres liquides, portefeuilles petits par rapport
au volume des titres, et des coûts bien approchés par une fourchette et une commission par unité échangée. Le momentum 12-1 mensuel sur les 500 titres les plus liquides
synthétiques échange 1.9~\% de son portefeuille par jour~; son ratio de Sharpe passe de 0.32 avant coûts à 0.21 après 10 points de base par unité échangée,
0.19 après l'emprunt, et 0.17 lorsqu'il est échangé à la clôture suivante (\cref{fig:rs:vectorised-backtests:waterfall}). Aucune de ces étapes ne dépend de la manière dont les
transactions sont exécutées, et un \omterm{def:rs:vectorised-backtests:backtest}{backtest} de niveau~2 ou~3 n'ajouterait que peu de chose. Il est faux, au sens où il manque la réponse, lorsque les exécutions dépendent du
carnet d'ordres (ordres passifs, chapitre~17), lorsque les transactions propres à la stratégie déplacent les prix (chapitre~27), lorsque le moment de la journée compte, et lorsque
des événements tels que les suspensions de cotation, les enchères ou les opérations sur titres interviennent entre la décision et la transaction.

\section{Les mensonges qu'il raconte}

\begin{proposition}[Coûts face à la rotation]\label{prop:rs:vectorised-backtests:breakeven}
Un portefeuille de rendement brut moyen $g$ par période et de rotation $\tau$ par période (moitié du poids échangé) atteint le point mort à un coût linéaire de
$g/(2\tau)$ par unité de poids échangé.
\end{proposition}

\begin{proof}
Le poids échangé par période vaut $2\tau$, et à un coût $c$ par unité, le coût par période vaut $2\tau c$~; on l'égalise à $g$.
\end{proof}

L'inversion à un jour sur le marché synthétique, reconstruite chaque jour comme un portefeuille neutre en dollars d'exposition brute~1, corrigée un mensonge à la fois
(\texttt{rs\_vecbt.waterfall}~; chaque étape conserve les corrections qui la précèdent)~:

\begin{center}
\small
\begin{tabular}{@{}p{5.4cm}ccc@{}}
\toprule
étape & ratio de Sharpe & rendement moyen annuel & rotation quotidienne\\
\midrule
naïve (survivants, exécution à la clôture du signal) & 7.17 & 27.0~\% & 74~\%\\
titres cotés à chaque décision & 7.22 & 27.4~\% & 74~\%\\
les 500 titres les plus liquides & 6.95 & 27.4~\% & 75~\%\\
10 pb par unité échangée & $-2.58$ & $-10.2~\%$ & 75~\%\\
emprunt à 50 pb par an & $-2.64$ & $-10.4~\%$ & 75~\%\\
exécution à la clôture suivante & $-9.45$ & $-37.7~\%$ & 75~\%\\
\bottomrule
\end{tabular}
\end{center}

\textbf{L'univers.} L'univers naïf est la liste des titres encore cotés à la fin de l'échantillon, celle que le chercheur télécharge aujourd'hui. Pour
l'inversion, il change peu (7.17 contre 7.22 sur la liste à date), pour le momentum beaucoup, et dans le sens qu'un manuel
ne prévoit pas~: 0.21 contre 0.44 (\cref{fig:rs:vectorised-backtests:waterfall}). Les titres radiés par la suite étaient les perdants du momentum, et la
vente à découvert profitait de leurs dernières baisses~; le \omterm{def:rs:point-in-time-data-and-the-biases:survivorship}{biais de survie} supprime ce profit. Son signe dépend de la stratégie~; sa taille ne peut qu'être mesurée.

\textbf{Coûts.} Le rendement brut de l'inversion est élevé et sa rotation plus élevée encore~: 74~\% du portefeuille par jour. À 10 points de base par unité de poids échangé
le coût vaut $2 \times 0.75 \times 0.001 \times 252 = 38~\%$ par an pour un brut de 27~\%, et le ratio de Sharpe passe de 6.95 à $-2.58$. Par
la \cref{prop:rs:vectorised-backtests:breakeven}, le portefeuille atteint le point mort à 7.3 points de base. Novy-Marx et Velikov (2016) ont constaté la même chose sur de vraies
anomalies~: la plupart des stratégies dont la rotation mensuelle est inférieure à 50~\% dégageaient encore des écarts significatifs après coûts lorsqu'elles étaient conçues pour les limiter, et peu
de celles à rotation plus élevée le faisaient.

\textbf{Emprunt.} Les positions courtes versent des frais au prêteur du titre (livre~1, chapitre~6), 50 points de base par an sur le collatéral général et bien davantage
sur les titres difficiles à emprunter~; ici, il coûte 0.25~\% par an et ne pèse guère sur le ratio de Sharpe. Ce point compte davantage pour les stratégies dont les positions courtes se concentrent sur de petits
titres encombrés, ce qui n'est pas le cas de celles de l'inversion.

\textbf{Calendrier.} Un portefeuille calculé à partir d'un cours de clôture ne peut pas être exécuté à ce cours~: les ordres de l'enchère de clôture sont transmis avant que la clôture ne soit connue
(livre~1, chapitre~13). Si elle est exécutée à la clôture suivante, l'inversion ne conserve plus aucun alpha (son information dure un jour, chapitre~13) et garde ses
coûts~: $-9.45$. La réponse honnête au niveau~1 utiliserait un signal calculé à partir des cours connus avant la clôture, et un modèle de coûts de l'enchère de clôture~;
le \omterm{def:rs:vectorised-backtests:backtest}{backtest vectorisé} peut retenir cette hypothèse, à condition que quelqu'un la consigne par écrit.

\textbf{Capitalisation.} Sommer les rendements quotidiens au lieu de les composer est sans conséquence pour de petits rendements, mais non pour les grands. Le capital
de l'inversion avec coûts, composé, tombe à 0.35 de sa valeur initiale au bout de dix ans~; sommé, il atteint $-0.04$, un capital qu'aucun portefeuille ne peut avoir.

\textbf{Anticipation.} Multiplier le portefeuille du jour par le rendement du jour fait perdre à l'inversion 694~\% par an ($-65$ de ratio de Sharpe)~: son signal est l'opposé de ce
rendement. L'anticipation ne flatte pas toujours~; elle fausse toujours.

\begin{omfigure}
\begin{tikzpicture}
\begin{axis}[name=rev,width=6.8cm,height=5.6cm,ybar,bar width=10pt,title={\small inversion à un jour},ylabel={ratio de Sharpe},tick label style={font=\footnotesize},
  label style={font=\small},xtick={0,1,2,3,4,5},xticklabels={naïf,PIT,liquide,coûts,emprunt,suivant},x tick label style={rotate=45,anchor=east},
  ymin=-10.5,ymax=8.5,grid=major,grid style={gray!25},table/col sep=comma]
\addplot[fill=omThm!45,draw=omThm] table[x=k,y=reversal]{figdata/research/16-vectorised-backtests/waterfall.csv};
\end{axis}
\begin{axis}[at={(rev.east)},anchor=west,xshift=1.6cm,width=6.8cm,height=5.6cm,ybar,bar width=10pt,title={\small momentum 12-1, mensuel},
  tick label style={font=\footnotesize},label style={font=\small},xtick={0,1,2,3,4,5},xticklabels={naïf,PIT,liquide,coûts,emprunt,suivant},
  x tick label style={rotate=45,anchor=east},ymin=0,ymax=0.5,grid=major,grid style={gray!25},table/col sep=comma]
\addplot[fill=omDef!55,draw=omDef] table[x=k,y=momentum]{figdata/research/16-vectorised-backtests/waterfall.csv};
\end{axis}
\end{tikzpicture}

\omcaption{Ratio de Sharpe de deux stratégies sur \texttt{firm.synthmkt} à mesure que les erreurs du \omterm{def:rs:vectorised-backtests:backtest}{backtest} sont corrigées une à une~: univers des survivants (naïf),
univers en temps réel (PIT), 500 titres les plus liquides, 10 pb par unité échangée, emprunt à 50 pb par an, exécution à la clôture suivante. Les deux
échelles diffèrent.}\label{fig:rs:vectorised-backtests:waterfall}
\end{omfigure}

\begin{omfigure}
\begin{tikzpicture}
\begin{axis}[width=11cm,height=5.6cm,xlabel={années},ylabel={capital ($\log_{10}$)},tick label style={font=\small},label style={font=\small},
  xmin=0,xmax=10,ymin=-2,ymax=1.5,grid=major,grid style={gray!25},legend columns=2,
  legend style={font=\footnotesize,at={(0.5,-0.25)},anchor=north,draw=none,fill=none,
  /tikz/every even column/.append style={column sep=8pt}},table/col sep=comma]
\addplot[black!50,thick,dashed] table[x=year,y=gross]{figdata/research/16-vectorised-backtests/capital.csv};
\addplot[omThm,thick] table[x=year,y=costed]{figdata/research/16-vectorised-backtests/capital.csv};
\addplot[omThm,thick,dotted] table[x=year,y=next]{figdata/research/16-vectorised-backtests/capital.csv};
\addplot[omDef,very thick] table[x=year,y=momentum]{figdata/research/16-vectorised-backtests/capital.csv};
\legend{inversion avant coûts,inversion après coûts et emprunt,inversion exécutée à la clôture suivante,momentum après toutes les corrections}
\end{axis}
\end{tikzpicture}

\omcaption{Capital composé de l'inversion sur les 500 titres synthétiques les plus liquides avant coûts, après coûts et emprunt, et exécuté à la clôture suivante,
avec le momentum mensuel après toutes les corrections. Données~: \texttt{rs\_vecbt} sur \texttt{firm.synthmkt}.}\label{fig:rs:vectorised-backtests:capital}
\end{omfigure}

\section{Une liste de contrôle}

Tout \omterm{def:rs:vectorised-backtests:backtest}{backtest} de niveau~1 que la société mène énonce, et \texttt{firm.vecbt} l'impose lorsque c'est possible~:
\begin{enumerate}
\item l'\omterm{def:rs:vectorised-backtests:timing}{instant de décision} de chaque donnée d'entrée, et le \omterm{def:rs:vectorised-backtests:timing}{délai d'exécution} entre celui-ci et la transaction (jamais 0~; 1 seulement avec une raison énoncée)~;
\item l'\omterm{def:rs:universe-symbology-and-corporate-actions:universe}{univers investissable} à chaque décision, en temps réel, avec les \omterm{def:rs:universe-symbology-and-corporate-actions:delisting}{rendements de radiation}~;
\item le coût par unité échangée, par titre si la liquidité diffère, et le coût de point mort (\cref{prop:rs:vectorised-backtests:breakeven})~;
\item les frais d'emprunt et la disponibilité des titres pour les positions courtes, ainsi que le financement de tout effet de levier et les intérêts sur la trésorerie~;
\item la \omterm{def:rs:vectorised-backtests:timing}{fréquence de rebalancement} et la dérive du portefeuille entre deux rebalancements~;
\item la capitalisation, et un capital qui ne peut pas descendre sous zéro~;
\item l'exposition brute, et les plafonds de position face à la liquidité~;
\item le nombre de variantes essayées avant celle-ci (chapitre~20), et le niveau auquel le résultat doit être confirmé.
\end{enumerate}
Arnott, Harvey et Markowitz (2019) ont exposé un tel protocole pour la recherche en apprentissage automatique, en soulignant que ses méthodes exigent bien plus de données
que n'en possède la finance.

\section{Tutoriel~: Sharpe 7.2 le lundi}\label{tut:rs:vectorised-backtests:tut}

\begin{tutorial}
\textbf{Objectif.} Backtester l'inversion à un jour et le momentum mensuel sur \texttt{firm.synthmkt} en un seul passage vectorisé chacun, corriger les erreurs une à une
et calculer le coût de point mort de l'inversion. \textbf{État final~:}
les \cref{fig:rs:vectorised-backtests:waterfall,fig:rs:vectorised-backtests:capital}~; le tableau.
\begin{tutsteps}
\item \textbf{La boucle principale}~: le portefeuille décalé, la transaction issue du portefeuille dérivé, les coûts, l'emprunt, le financement.
\omcode{code/firm/vecbt/firm_vecbt.py}{90}{101}{La boucle de période du backtest de niveau~1.}\label{lst:rs:vectorised-backtests:loop}
\item \textbf{L'\omterm{def:rs:universe-symbology-and-corporate-actions:universe}{univers investissable}}~: les titres les plus liquides selon le volume en dollars glissant, tel qu'il est connu à chaque clôture.
\omcode{code/research/16-vectorised-backtests/python/rs_vecbt.py}{37}{43}{Un univers liquide en temps réel.}\label{lst:rs:vectorised-backtests:liquid}
\item \textbf{Exécuter} \texttt{waterfall()} pour les deux signaux, \texttt{lookahead()}, \texttt{breakeven\_cost()}, \texttt{compounding()} et
\texttt{fig\_vecbt.py}.
\end{tutsteps}
\textbf{Pour aller plus loin.} Lisser le signal d'inversion sur trois jours et observer la rotation et le coût de point mort~; facturer les coûts par titre, proportionnellement
à la volatilité de chaque titre.
\end{tutorial}

\section{Construction~: le backtester de niveau~1}\label{bld:rs:vectorised-backtests:vecbt}

\begin{build}
\bfield{Rôle} Le backtester de criblage de la société~: quelques secondes par exécution, honnête par construction pour le calendrier, l'univers, les coûts et le financement, et source
du type \texttt{BacktestResult} que tous les \omterm{def:rs:vectorised-backtests:backtest}{backtests} ultérieurs renvoient.
\bfield{Interface} \texttt{BacktestResult(dates, names, weights, trades, gross, costs, net, capital, meta)} avec \texttt{turnover},
\texttt{gross\_exposure}, \texttt{net\_exposure}~; \texttt{signal\_to\_weights(signal, universe, gross, neutral)}~; \texttt{backtest(weights, returns, lag,
universe, cost, borrow, cash\_rate, borrow\_spread, cap, compound, periods)}.
\bfield{Règles} Le décalage est un argument, non une hypothèse~; l'univers s'applique à l'\omterm{def:rs:vectorised-backtests:timing}{instant de décision}~; les transactions se mesurent à partir du portefeuille dérivé~; les coûts
sont imputés à chaque transaction~; le capital se compose par défaut.
\bfield{Tests d'acceptation} \texttt{code/firm/vecbt/tests/}~: décalages 0, 1 et 2 sur une série manuelle~; capitalisation comparée à la sommation~; la transaction issue d'un portefeuille dérivé~;
coûts et rotation d'un portefeuille alternant~; emprunt, intérêts sur la trésorerie et financement par levier~; masques d'univers et plafonds~; poids neutres en dollars d'une exposition brute donnée.
\bfield{Pour aller plus loin} Coûts issus d'un modèle d'impact en racine carrée (chapitre~27)~; disponibilité des titres à emprunter par titre~; trésorerie multidevise.
\end{build}

\begin{omsources}
\item R.~Novy-Marx et M.~Velikov, ``A taxonomy of anomalies and their trading costs'', \emph{Review of Financial Studies} 29(1), 2016.
\item R.~Arnott, C.~R.~Harvey et H.~Markowitz, ``A backtesting protocol in the era of machine learning'', \emph{Journal of Financial Data Science} 1(1), 2019.
\end{omsources}

\section{Exercices}

\begin{exercise}[$\star$]\label{exo:rs:vectorised-backtests:1}
Un portefeuille rapporte 12~\% brut par an et tourne de 20~\% de lui-même chaque jour (252 jours). À quel coût par unité échangée atteint-il le point mort~?
\end{exercise}

\begin{exercise}[$\star$]\label{exo:rs:vectorised-backtests:2}
Un portefeuille détient 60~\% en A et 40~\% en B. A rapporte 10~\% et B 0~\% sur la période. Quels sont les poids dérivés, et que faut-il échanger pour revenir à
60/40~?
\end{exercise}

\begin{exercise}[$\star$]\label{exo:rs:vectorised-backtests:3}
Les rendements quotidiens de $+50~\%$ et $-40~\%$ alternent pendant dix jours. Quel est le capital après dix jours en composant les rendements, et que suggère la somme des rendements~?
\end{exercise}

\begin{exercise}[$\star\star$]\label{exo:rs:vectorised-backtests:4}
Pourquoi le \omterm{def:rs:vectorised-backtests:backtest}{backtest} avec anticipation (look-ahead) du retournement perd-il de l'argent, alors que celui d'un signal de momentum incluant le rendement du jour en gagnerait~?
\end{exercise}

\begin{exercise}[$\star\star$]\label{exo:rs:vectorised-backtests:5}
Pourquoi le \omterm{def:rs:point-in-time-data-and-the-biases:survivorship}{biais de survie} peut-il gonfler le \omterm{def:rs:vectorised-backtests:backtest}{backtest} d'une stratégie et réduire celui d'une autre~? Donner le mécanisme pour le momentum sur le marché synthétique.
\end{exercise}

\begin{exercise}[$\star\star$]\label{exo:rs:vectorised-backtests:6}
Un portefeuille neutre en dollars, d'exposition brute 2, dépose son capital en collatéral rémunéré à 3~\% par an et verse 50 points de base par an pour emprunter ses positions courtes. Quel
est l'effet du financement, gain ou coût, par an et par unité de capital~?
\end{exercise}

\begin{exercise}[$\star\star\star$]\label{exo:rs:vectorised-backtests:7}
\emph{Programmation.} Reconstruire le retournement à partir de la moyenne des rendements des trois derniers jours au lieu du seul dernier. Comment la \omterm{def:rs:vectorised-backtests:turnover}{rotation du portefeuille} (turnover), le ratio de Sharpe brut et
le coût de point mort évoluent-ils~?
\end{exercise}

\begin{exercise}[$\star\star\star$]\label{exo:rs:vectorised-backtests:8}
\emph{Trouver la faille.} «~Notre stratégie s'exécute à la clôture au cours de clôture, ce qui est correct parce que nous utilisons des ordres au marché à la clôture.~»
\end{exercise}

\section{Problème~: Sharpe 7.2 le lundi, $-2.6$ le vendredi}

\begin{problem}[{Problème du week-end --- un backtest, mensonge par mensonge}]\label{pb:rs:vectorised-backtests:1}
Le retournement à un jour et le momentum mensuel 12-1 sur \texttt{firm.synthmkt}, neutre en dollars, d'exposition brute 1.

\medskip\noindent\textbf{Partie I --- Le \omterm{def:rs:vectorised-backtests:backtest}{backtest} naïf.}
\begin{enumerate}
\item Quel ratio de Sharpe et quel rendement moyen le \omterm{def:rs:vectorised-backtests:backtest}{backtest} naïf du retournement affiche-t-il, et que suppose-t-il~?
\item Que montre la version à retard nul, et pourquoi son signe est-il négatif~?
\item Quelle \omterm{def:rs:vectorised-backtests:turnover}{rotation du portefeuille} (turnover) présente le retournement~?
\item Quelles entrées du \omterm{def:rs:vectorised-backtests:backtest}{backtest} naïf un chercheur aurait-il pu connaître à l'\omterm{def:rs:vectorised-backtests:timing}{instant de décision}~?
\end{enumerate}

\medskip\noindent\textbf{Partie II --- L'univers.}
\begin{enumerate}[resume]
\item Que font les univers point-in-time et liquide à chacune des stratégies~?
\item Pourquoi le \omterm{def:rs:point-in-time-data-and-the-biases:survivorship}{biais de survie} abaisse-t-il le momentum ici~?
\item Que ferait-il à une stratégie qui achète des actions bon marché et en difficulté~?
\end{enumerate}

\medskip\noindent\textbf{Partie III --- Coûts et délais.}
\begin{enumerate}[resume]
\item Que font 10 points de base par unité échangée au retournement et au momentum~?
\item Quel est le coût de point mort du retournement~?
\item Combien coûte l'emprunt~?
\item Que se passe-t-il lorsque le retournement est exécuté à la clôture suivante, et pourquoi~?
\item Qu'est-ce que la composition des rendements change pour le retournement après coûts~?
\end{enumerate}

\medskip\noindent\textbf{Partie IV --- Le verdict.}
\begin{enumerate}[resume]
\item Énoncer le \emph{résultat nommé}~: la cascade du ratio de Sharpe, du \omterm{def:rs:vectorised-backtests:backtest}{backtest} naïf au \omterm{def:rs:vectorised-backtests:backtest}{backtest} honnête de niveau 1, mensonge par mensonge, pour les deux stratégies.
\item Laquelle des deux stratégies relève d'un résultat de niveau 1, et laquelle exige un \omterm{def:rs:vectorised-backtests:fidelity}{niveau de fidélité} supérieur avant que quiconque y croie~?
\item Que changeraient un modèle de coût de l'enchère de clôture et un signal fondé sur les cours de 15:50~?
\item Comment réduire la \omterm{def:rs:vectorised-backtests:turnover}{rotation du portefeuille} (turnover) du retournement sans perdre son alpha~?
\item Quelle ligne de la liste de contrôle le \omterm{def:rs:vectorised-backtests:backtest}{backtest} naïf a-t-il violée en premier~?
\item Que faut-il attendre, d'après les résultats de Novy-Marx et Velikov, d'un véritable retournement quotidien~?
\item Que doit consigner le \omterm{def:rs:the-research-process:log}{journal de recherche} sur ce \omterm{def:rs:vectorised-backtests:backtest}{backtest}~?
\item En une phrase~: à quoi sert un \omterm{def:rs:vectorised-backtests:backtest}{backtest vectorisé}~?
\end{enumerate}
\end{problem}

\section{Questions d'entretien}

\begin{interviewq}[$\star$ \iqroles{chercheur}]\label{iq:rs:vectorised-backtests:1}
Citez cinq façons dont un \omterm{def:rs:vectorised-backtests:backtest}{backtest} peut surestimer la performance.
\end{interviewq}

\begin{interviewq}[$\star\star$ \iqroles{chercheur, développeur}]\label{iq:rs:vectorised-backtests:2}
Décrivez, en mots, les étapes d'un \omterm{def:rs:vectorised-backtests:backtest}{backtest vectorisé} d'une stratégie long-short quotidienne, et l'endroit où chaque biais intervient.
\end{interviewq}

\begin{interviewq}[$\star\star$ \iqroles{chercheur, trader}]\label{iq:rs:vectorised-backtests:3}
Le ratio de Sharpe brut d'une stratégie vaut 7. Pourquoi pourrait-on encore ne pas la mettre en production~?
\end{interviewq}

\begin{interviewq}[$\star\star$ \iqroles{chercheur}]\label{iq:rs:vectorised-backtests:4}
Quand un \omterm{def:rs:vectorised-backtests:backtest}{backtest vectorisé} est-il suffisant, et quand faut-il un \omterm{def:rs:vectorised-backtests:backtest}{backtest} événementiel~?
\end{interviewq}

\begin{interviewq}[$\star\star$ \iqroles{chercheur, risque}]\label{iq:rs:vectorised-backtests:5}
Comment tenez-vous compte de l'emprunt et du financement dans un \omterm{def:rs:vectorised-backtests:backtest}{backtest} long-short~?
\end{interviewq}

\begin{interviewq}[$\star\star\star$ \iqroles{chercheur}]\label{iq:rs:vectorised-backtests:6}
Dérivez le coût de point mort d'une stratégie à partir de son rendement brut et de son turnover, et utilisez-le pour comparer un retournement quotidien avec un momentum mensuel.
\end{interviewq}
